\section{Related Work}

\subsection{Preload Monitoring in Bolted Joints}

Direct preload monitoring has been addressed using contact sensors and wave-based approaches. Strain-gauge bolts and load cells provide direct force-related measurements, while ultrasonic and coda-wave techniques use wave propagation changes to infer preload or preload loss \cite{chen2022coda,chen2023fullrange}. Fiber Bragg grating sensors and other embedded sensing concepts have also been proposed for connection-state monitoring \cite{deng2022fbg}. These approaches can be accurate, but they require additional sensing hardware, signal-conditioning infrastructure, and calibration. The present study does not aim to replace these methods in high-accuracy applications; rather, it investigates whether an optical approach can provide a scalable rough estimate of preload for standard fasteners under controlled conditions.

\subsection{Bolt-Head Strain as a Preload Indicator}

The mechanical basis of the proposed framework is supported by previous work showing that the bolt-head surface contains preload-sensitive strain information. Wang et al. investigated the strain distribution on the bolt-head top surface and showed that radial strain has an approximately linear relationship with pretightening force \cite{wang2020bolt}. Their results indicate that the bolt-head face is not only a convenient mounting location for a strain gauge, but also a mechanically informative surface. Wang et al. later proposed a wireless bolt-head strain sensor with self-temperature compensation, using the different sensitivities of radial and circumferential strain regions on the bolt-head face \cite{wang2024head}. Ye et al. also reported a bolt-force monitoring method based on radial strain on the bolt-head surface in maglev infrastructure maintenance, with strain gauges arranged on a circular ring to reduce the influence of lateral forces and bending moments \cite{ye2024maglev}. These studies provide an important premise for PreForceVision: preload produces a spatially structured response on the visible bolt-head face.

\subsection{Digital Image Correlation and Optical Deformation Measurement}

Digital image correlation (DIC) is a full-field optical method for estimating displacement and strain from reference and deformed images. Pan et al. provide a comprehensive review of two-dimensional DIC, including subset matching, displacement estimation, strain computation, and accuracy considerations \cite{pan2009dic}. DIC has also been used in high-accuracy optical extensometry and mechanical testing \cite{zhu2017extensometer}. In the bolt context, Huo et al. proposed a DIC-based bolt preload monitoring method in which an industrial camera records speckle images of the bolt head before and after deformation, and the resulting surface strain field is related to preload \cite{huo2026dic}. Li et al. similarly proposed early bolt-loosening detection using DIC on a bolt-head surface to predict axial force from average strain \cite{li2024dic}. These studies are close to the present work, but they rely on controlled conditions and intentionally textured surfaces. PreForceVision uses DIC not only as a measurement method, but also as a physically interpretable baseline and supervision source for a later image-based predictive model.

\subsection{Vision-Based Bolt Inspection and Deep Learning}

Vision-based bolt monitoring has largely focused on visible loosening, rotation, or damage. Early work combined Hough-transform features and support vector machines for loosened-bolt detection \cite{cha2016vision}. Later studies used deep detectors and synthetic data to identify bolts and bolt-loosening states \cite{pham2020bolt,ta2020monitoring,ta2022corroded,sun2022yolo, afifi2026gaussian}. Other work used keypoint detection, angle estimation, or geometric imaging theory to quantify loosening angles or exposed-bolt lengths \cite{gong2022quantitative,qian2024loosening}. These approaches demonstrate the feasibility of image-based inspection under varying distances, angles, and lighting, but their target is generally a visible geometric state, not preload. Recent review literature confirms that vision-based deep learning for bolted joints is dominated by damage, looseness, and detection tasks rather than quantitative preload estimation \cite{malik2024review}.

Deep learning has also been applied to optical deformation measurement\cite{afifi2025data}. Boukhtache et al. introduced a learning-based approach that connects deep networks with DIC-like displacement estimation \cite{boukhtache2021strainnet}. Yang et al. proposed Deep DIC for end-to-end displacement and strain estimation \cite{yang2022deepdic}. These studies suggest that learned models can accelerate or approximate deformation-field estimation, but they do not remove the need for careful physics-based validation. In the present application, the signal is small, the surface is metallic, and lighting or registration errors can easily dominate the apparent difference between two images. The learning problem must therefore be constrained by mechanics and evaluated against ground truth from finite element analysis, strain gauges, optical 3D measurements, and DIC.

%%%%%%%%%%%%%%%%%%5
\subsection{Research Gap and Question}
The research gap lies in the missing link between surface-accessible optical measurement and mechanically meaningful internal state estimation. Sensor-based preload methods can provide accurate force-related measurements, but they require modified fasteners, additional hardware, wiring, coupling, or calibration. Vision-based bolt inspection is scalable, but it mainly detects visible states such as missing bolts, corrosion, rotation, or looseness, rather than preload as the functional quantity of the joint. DIC-based methods provide a stronger physical basis, but they still require controlled surface preparation, stable acquisition, and independent validation.

This paper therefore investigates whether the preload-induced deformation of the visible bolt-head front face can be used as an optical signature for rough preload estimation. The research question is: How can small preload-induced surface deformations be extracted and validated systematically so that they support controlled, uncertainty-aware preload estimation?
%The research gap lies between accurate but hardware-dependent preload measurement and scalable but mechanically indirect visual inspection. Sensor-based methods can measure preload-related quantities, but require additional components or instrumentation. Vision-based bolt inspection can scale to many fasteners, but most methods detect looseness states, rotation, markings, or damage rather than preload. DIC-based methods provide a stronger physical foundation, but still require controlled acquisition, surface texture, and careful calibration.

%The research question is therefore: How can we systematically extract the small preload-induced deformation of the bolt-head front face with sufficient repeatability to support rough preload estimation under controlled conditions?
